\scanpage{29}
\mspage{9}
\deleted{\unsure{entièrement déterminée par}}
\inserted{\uline{équivalente à}}
\uline{la connaissance de}
\figure{\begin{tikzpicture}[baseline=(t.base)]
\node[draw,ellipse,inner sep=3pt] (t) {\uline{tous}};
\end{tikzpicture}}
\uline{les} $H^i(X,\mathbb Z_\ell)$.
De façon plus précise, pour un $\ell$ donné, la connaissance de $H^i(X,\mathbb Z_\ell)$ équivaut à celle du rang \inserted{$b_i(X)$} de $H^i(X^{\mathrm{an}},\mathbb Z)$, \inserted{c'est à dire de $H^i(X,\mathbb Z_\ell)$ sur $\mathbb Z_\ell$,} et de la composante $\ell$-primaire du sous-groupe de torsion de $H^i(X^{\mathrm{an}},\mathbb Z)$.

On obtient, en passant, par voie transcendante, le fait que si le corps de base $k$ est de car. $0$, alors les \inserted{$H^i(X,\mathbb Z_\ell)$ sont de t.f., leur rang,} $b_i(X,\mathbb Z_\ell)$ \deleted{est} pour $\ell$ variable est indépendant de $\ell$, et que pour $X$ fixé, seul$^{\mathrm t}$ pour un nb fini de $\ell$ \unsure{le} $\unsure{H^*|H^i}(X,\mathbb Z_\ell)$ a de la torsion.
Et par la $\overline{\mathrm m}$ voie, on trouve \inserted{p.ex.} que si on a un endom. $f$ de $X$, alors pour $i$ fixé, \deleted{l'ensemble} le polynôme caractéristique de l'endom. de $H^i(X,\mathbb Z_\ell)$
\margin{28}
